\chapter{Experimental Results and Discussion}\label{chap:9}
\label{chap:experiments}

\noindent This chapter discusses the experimental work conducted in developing the NIDPS.

\section{Experimental Setup}\label{sec:9.1}

\noindent All the experiments discussed in this chapter were performed in a small isolated test network. The NIDPS device is the single Linux host running hostapd as the wireless access point (Section~\ref{sec:overview}), Suricata, the ARP monitor and the Layer 2 anomaly service. The attacking machine was a laptop which was used to run the scripts. A second laptop was plugged through the access point and so was a phone, both of which were used as victims. For each of the attack scripts listed in Table~\ref{tab:attackscripts}, the following steps were taken:
\begin{enumerate}
    \item The ARP monitor, Suricata, and the Layer 2 anomaly service were started for a brief warm-up time so that the host-behaviour tracker of Section~\ref{sec:7.1.7.2} had a head start.
    \item The attack script was executed once on the attacker machine, via the attack runner described in Section~\ref{sec:7.1.8.1}, which records the attack name, start and end time, and the attacker and victim addresses.
    \item The Suricata eve.json log, the evidence table that the services wrote, and the anomaly service's log were examined by hand after the run to determine from which detector any evidence came, and whether the attack produced any at all.
\end{enumerate}

\begin{table}[H]
\centering
\caption{Attack scripts used in the experiments (Section~\ref{sec:7.1.8.1}).}\label{tab:attackscripts}
\begin{tabular}{|>{\raggedright\arraybackslash}p{0.32\linewidth}|>{\raggedright\arraybackslash}p{0.60\linewidth}|}
\hline
\textbf{Attack} & \textbf{Tool / script} \\ \hline
Port scan & nmap \\ \hline
SSH brute force & hydra \\ \hline
SYN flood & hping3 \\ \hline
DDoS (single source) & hping3, run from one attacker machine \\ \hline
ARP spoofing & arpspoof \\ \hline
DNS spoofing & dnsspoof (run together with ARP spoofing of the same victim) \\ \hline
C2 beaconing & custom Python TLS client \\ \hline
C2 beaconing with a chosen JA3 & custom Python TLS client \\ \hline
Data exfiltration & custom Python TLS client \\ \hline
Lateral movement & custom Python TLS client \\ \hline
\end{tabular}
\end{table}

\noindent The Isolation Forest was used with scikit-learn on a benign capture from the test network, with the same 78 features that the live extractor generates. The flows were then scored using the scoring and calibration code from Section~\ref{sec:7.1.7.1}.

\section{Signature-Based Detection: Suricata Against the Self-Collected Attacks}\label{sec:9.2}

\noindent Only four custom Suricata rules are enabled: SSH brute force, port scan, SYN flood and DNS spoofing (Section~\ref{sec:7.1.6.1}). The below was the result of ten attacks run on the live system.

\noindent Six of the ten attacks generated no Suricata alert at all, and a seventh, DDoS, was detected only some of the time. This outcome is in accordance with the prediction in the approach chapter: signature detection only handles attacks with a consistent pattern at the packet level, and the remaining attacks were selected precisely because they don't have one. This is the exception of ARP spoofing. While Suricata itself did not catch it, the ARP monitor did catch it, which is the reason for the detector being a separate component rather than another Suricata rule.

\begin{table}[H]
\centering
\caption{Suricata's response to each self-collected attack.}\label{tab:suricataresults}
\begin{tabular}{|>{\raggedright\arraybackslash}p{0.26\linewidth}|>{\raggedright\arraybackslash}p{0.14\linewidth}|>{\raggedright\arraybackslash}p{0.54\linewidth}|}
\hline
\textbf{Attack} & \textbf{Alert raised} & \textbf{Notes} \\ \hline
Port scan & Yes & Matched the twenty-ports-in-10\,s rule every time it was run. \\ \hline
SSH brute force & Yes & Matched the five-SYNs-to-port-22 rule as soon as the fifth connection attempt landed inside the 60\,s window. \\ \hline
SYN flood & Yes & Matched the 100-SYNs-in-10\,s rule. \\ \hline
DNS spoofing & Yes & Matched because the forged reply came from an address other than the lab resolver. \\ \hline
DDoS (single source) & Partially & Crossed the SYN-flood threshold when the flood rate was high enough, but a slower run of the same script stayed under 100 SYNs in 10\,s and produced no alert. A genuine distributed flood, coming from many sources at once, would likely never trip this single-source rule at all. \\ \hline
ARP spoofing & No & No Suricata rule inspects ARP traffic; this attack was instead picked up by the ARP monitor (Section~\ref{sec:7.1.6.3}), which raised evidence for the spoofed address as designed. \\ \hline
C2 beaconing & No & No packet-level signature exists for a client that just repeats regular TLS connections to one address; this is the behaviour Layer 2 is meant to catch. \\ \hline
C2 beaconing with a chosen JA3 & No & The JA3-blacklist rule is written but left commented out (Section~\ref{sec:7.1.6.1}) because no bad JA3 hash had been recorded for the lab yet, so the attack produced no Suricata alert. \\ \hline
Data exfiltration & No & A single large upload to a new destination has no signature; this is exactly the behaviour the exfiltration check in Section~\ref{sec:7.1.7.2} is designed to flag instead. \\ \hline
Lateral movement & No & Internal connections between lab hosts matched no active rule. \\ \hline
\end{tabular}
\end{table}


\section{Anomaly Detection: Isolation Forest Against the Self-Collected Attacks}\label{sec:9.3}

\noindent The four attacks that did not leave any trace in the Suricata log but are still recorded at the flow level are: C2 beaconing, C2 beaconing with a chosen JA3, data exfiltration and lateral movement. The scores for these were determined with the help of the provisional Isolation Forest and host-behaviour checks described in Section~\ref{sec:9.1}. The evidence captured by the anomaly service, and the flows scored from the warm-up capture, were summarised and listed in Table~\ref{tab:ifresults} below for both sides, although with a small sample for each.

\begin{table}[H]
\centering
\caption{Isolation Forest and host-behaviour results on the four attacks Suricata missed.}\label{tab:ifresults}
\begin{tabular}{|>{\raggedright\arraybackslash}p{0.28\linewidth}|>{\raggedright\arraybackslash}p{0.68\linewidth}|}
\hline
\textbf{Attack} & \textbf{Observation} \\ \hline
Data exfiltration & The outbound-bytes burst exceeded the host's rolling baseline and tripped the exfiltration check of Section~\ref{sec:7.1.7.2} on every run; this check does not depend on the Isolation Forest model itself. \\ \hline
C2 beaconing (fixed JA3) & Once five or more connections to the same destination and JA3 had been seen, the beaconing check flagged the host, because the gaps between connections and the bytes per connection stayed regular, exactly the pattern the coefficient-of-variation test is built to catch. \\ \hline
C2 beaconing with a chosen JA3 & Behaved the same way as the plain beaconing run; changing the JA3 string did not by itself change the regularity of the timing or byte pattern, so it was still flagged. \\ \hline
Lateral movement & The flows involved did not fit the exfiltration or beaconing patterns. The Isolation Forest scored a few of these flows above the ad hoc threshold used for this test, but not consistently, and several lateral-movement flows scored close to the benign flows from the warm-up capture. \\ \hline
\end{tabular}
\end{table}

\noindent Here, precision, recall and the false positive rate are not reported as a meaningful value would require many more flows and the cross-dataset comparison of Section~\ref{sec:crossdataset}, neither of which this round of testing attempted. Lateral movement yielded the least promising result and will probably need the Layer 3 classifier, when it becomes available, not the Isolation Forest alone to detect successfully.

\section{Discussion}\label{sec:9.4}

\noindent The division of labour between all three layers was tested in the following experiments. All the attacks that had a fixed packet-level pattern were successfully blocked by Suricata, and all the attacks that did not have a fixed packet-level pattern were not (such as JA3-blacklist, where in principle there was a signature, but it was not populated). Hence the ARP monitor is able to detect the attack type separately as an attack of this type. The two behaviours for which they were created were learned by the Isolation Forest and the host-behaviour checks, respectively. Lateral movement was weaker and is probably more suitable for the Layer 3 classifier, once it has been built.


\section{Limitations of the Current Experimental Work}\label{sec:9.5}

\begin{itemize}
    \item No measure was possible for the system performance measures: forwarding latency, throughput, time to block and recovery behaviour, because the enforcer, the watchdog and the Layer 3 classifier have not been built yet (Section~\ref{sec:7.1.2}).
    \item Each attack was executed a relatively small number of times. This is enough to observe whether a detector fires but not enough to create a stable false positive or detection rate.
    \item The DDoS script employed here is a single-source flood rather than a genuinely distributed attack, therefore a partial detection by the SYN-flood rule does not necessarily constitute a representative reflection of how that rule would perform against a real distributed attack.
\end{itemize}

\section{Conclusion}\label{sec:9.6}

\noindent A handful of experiments were shown in this chapter on the components of the NIDPS which are already implemented. Suricata detected the attacks which it has rules for and missed the rest, and the ARP monitor detected the one attack Suricata can't structurally see. The host-behaviour checks and the Isolation Forest achieved good results on exfiltration and beaconing and performed poorly on lateral movement. These results together offer an early indication that the layered design is working as intended.
